\section{Formal results}\label{app:formal}

The main paper states one result, \cref{thm:population}, and uses ten more. They are stated
here in the order the paper relies on them; proofs are in \cref{app:proofs}. Equation and
section references to the main paper resolve to its numbering, since both are compiled as one
document.

\subsection{What the group-level part measures}
\begin{proposition}[Composition of the transported gain]\label{prop:transported}
Under the quadratic score, write $\bar q_m(c)=\mathbb E[q_m(\cdot\mid X)\mid C=c]$ for
model $m$'s mean forecast in cell $c$, $\operatorname{Var}(q_m\mid c)$ for the
within-cell covariance matrix of its forecast vector, and $p(c)$ for the cell's transport
law. Then
\begin{equation}
\Delta P_c=
\underbrace{\Bigl[\lVert p(c)-\bar q_0(c)\rVert^2-\lVert p(c)-\bar q_1(c)\rVert^2\Bigr]}
_{\text{improvement in mean-forecast fit}}
\;-\;
\underbrace{\Bigl[\operatorname{tr}\operatorname{Var}(q_1\mid c)-
\operatorname{tr}\operatorname{Var}(q_0\mid c)\Bigr]}_{\text{increase in within-cell dispersion}} .
\label{eq:transported-composition}
\end{equation}
\end{proposition}

\noindent Only the first bracket is prior fit. The second rewards a model for making
\emph{less} dispersed predictions inside a cell at a fixed mean forecast, so $\dP$ can absorb
an entire observed gain with both models' mean forecasts sitting exactly on the cell
frequencies; \cref{app:transported} gives a four-point example. This is why the main paper
calls $\dP$ group-level rather than prior fit, and why confidence matching moves $\hdP$ so
sharply: a temperature change is almost purely a change in dispersion.

\subsection{Every Bregman score}
\begin{theorem}[Bregman-score covariance identity]\label{thm:bregman}
Let $H_x(y)=S_\psi(q_1(\cdot\mid x),y)-S_\psi(q_0(\cdot\mid x),y)$ for a Bregman score generated by a strictly convex, differentiable $\psi$, and write $\Delta g_y(x)=\partial_y\psi(q_1(\cdot\mid x))-\partial_y\psi(q_0(\cdot\mid x))$ for the contrast of the $y$-th gradient coordinate. Then
\begin{equation}
\Delta R_c=\sum_{y=1}^K\operatorname{Cov}\!\bigl(\Delta g_y(X),\,\ind\{Y=y\}\mid C=c\bigr).
\label{eq:bregman-cov-id}
\end{equation}
\end{theorem}

\begin{corollary}[Quadratic and log score]\label{cor:bregman}
For $\psi(q)=\lVert q\rVert_2^2$, $\Delta g_y(x)=2\Delta q_y(x)$ and Eq.~\eqref{eq:bregman-cov-id} is Eq.~\eqref{eq:cov-id}. For $\psi(q)=\sum_kq_k\log q_k$, $\Delta g_y(x)=\log q_1(y\mid x)-\log q_0(y\mid x)$, and
\begin{equation}
\Delta R_c=\sum_{y=1}^K\operatorname{Cov}\!\left(\log\frac{q_1(y\mid X)}{q_0(y\mid X)},\ \ind\{Y=y\}\ \middle|\ C=c\right).
\label{eq:log-cov-id}
\end{equation}
\end{corollary}

\noindent The log-score rows of \cref{tab:sggaudit} therefore carry the same exact covariance
meaning as the quadratic ones. The implemented log score floors each probability at
$\epsilon=10^{-6}$; the identity is exact for the unfloored score and negligibly different
whenever no relevant probability is within $\epsilon$ of zero.

\subsection{Coarsening, and a variable left out of the grouping}
\begin{proposition}[Coarsening decomposition]\label{prop:coarsen}
Let $\bar\phi=g(\phi)$ be any coarsening of the grouping variable, so that each cell
$\bar c$ of $\bar\phi$ is a union of cells of $\phi$. Then
\begin{equation}
\Delta R_{\bar c}
=\mathbb E\!\left[\Delta R_C \mid \bar\phi=\bar c\right]
+\sum_{y=1}^K\operatorname{Cov}\!\left(\mathbb E[H_X(y)\mid C],\,
\Pr(Y=y\mid C)\ \middle|\ \bar\phi=\bar c\right).
\label{eq:coarsen}
\end{equation}
\end{proposition}

\begin{proposition}[Sensitivity bound for an unrecorded confounder]\label{prop:sensitivity}
Suppose $H_x(y)\in[-M,M]$ for all $x,y$ (WAGER already requires a bounded score). Fix a
cell $c$ of $\phi$, write $a_y=\mathbb E[H_X(y)\mid\phi,Z]$, $b_y=\Pr(Y=y\mid\phi,Z)$
(both conditional on $\phi=c$) for the two $K$-vectors on the right of
Eq.~\eqref{eq:coarsen}, and let
\begin{equation}
\rho=\frac{\sum_{y=1}^K\operatorname{Cov}(a_y,b_y\mid c)}
{\sqrt{\sum_{y=1}^K\operatorname{Var}(a_y\mid c)}\sqrt{\sum_{y=1}^K\operatorname{Var}(b_y\mid c)}}
\in[-1,1]
\end{equation}
be their aggregate correlation. Then
\begin{align}
\left|\Delta R_c(\phi)-\mathbb E[\Delta R_{(\phi,Z)}\mid\phi=c]\right|
&\ \le\ |\rho|\sqrt{\textstyle\sum_{y=1}^K\operatorname{Var}(a_y\mid c)}\,
\sqrt{\textstyle\sum_{y=1}^K\operatorname{Var}(b_y\mid c)}\notag\\
&\ \le\ |\rho|\sqrt{\textstyle\sum_{y=1}^K\operatorname{Var}(H_X(y)\mid\phi=c)}\,
\sqrt{\textstyle\sum_{y=1}^Kp_y(c)\bigl(1-p_y(c)\bigr)},
\label{eq:sensitivity-bound}
\end{align}
where $p_y(c)=\Pr(Y=y\mid\phi=c)$; the second inequality holds termwise
($\operatorname{Var}(a_y\mid c)\le\operatorname{Var}(H_X(y)\mid c)$ and
$\operatorname{Var}(b_y\mid c)\le p_y(c)(1-p_y(c))$ by the law of total variance, since
conditioning on more variables cannot increase the variance of a conditional
expectation) followed by monotonicity of the square root applied to each sum
separately -- a product of two square-root sums, not a sum of per-label square roots,
which the law of total variance alone would not license.
\end{proposition}

\begin{corollary}[Free bound]\label{cor:sensitivity-crude}
Each $a_y$ ranges over $[-M,M]$ as a conditional expectation of $H$, so
$\operatorname{Var}(a_y\mid c)\le M^2$ and $\sum_y\operatorname{Var}(a_y\mid c)\le KM^2$.
The label probabilities satisfy $\sum_yb_y=1$, which gives $\sum_yb_y^2\le1$ and, by
Cauchy--Schwarz, $\sum_y(\mathbb E b_y)^2\ge1/K$; hence
\begin{equation}
\sum_{y=1}^K\operatorname{Var}(b_y\mid c)
=\mathbb E\Bigl[\textstyle\sum_yb_y^2\Bigr]-\sum_y\bigl(\mathbb E b_y\bigr)^2
\ \le\ 1-\frac1K .
\label{eq:simplex-var}
\end{equation}
Substituting both into Eq.~\eqref{eq:sensitivity-bound} gives the data-free bound
\begin{equation}
\left|\Delta R_c(\phi)-\mathbb E[\Delta R_{(\phi,Z)}\mid\phi=c]\right|
\ \le\ |\rho|\,M\sqrt{K-1}.
\label{eq:sensitivity-crude}
\end{equation}
\end{corollary}

\subsection{The estimator}
\begin{theorem}[Finite-sample WAGER identity]\label{thm:sample}
For every cell with $n_c\ge2$, without a distributional assumption,
\begin{equation}
\widehat{\Delta T}_c=\widehat{\Delta P}_c+\widehat{\Delta R}_c.
\end{equation}
Moreover, when the cell's examples are drawn i.i.d.\ from the cell population,
$\widehat{\Delta R}_c$ is an unbiased order-two U-statistic for $\Delta R_c$.
Independence, not exchangeability alone, is what unbiasedness requires: the
independent-coupling definition of $\Delta P_c$ pairs an example with a label drawn
independently from the same cell, so $\mathbb E[H_1(Y_2)]$ matches it only for
independent within-cell draws.
\end{theorem}

\begin{proposition}[Exact attenuation of the in-sample plug-in]\label{prop:attenuate}
For every cell with $n_c\ge2$ and every $i\in c$, writing $\hat P_i$ for
Eq.~\eqref{eq:linear} and $\hat R_i=H_i(y_i)-\hat P_i$,
\begin{equation}
\tilde P_i=\hat P_i+\frac1{n_c}\hat R_i,
\qquad
\tilde R_i:=H_i(y_i)-\tilde P_i=\frac{n_c-1}{n_c}\,\hat R_i.
\label{eq:attenuation}
\end{equation}
Consequently the in-sample plug-in's dataset-level statistic obeys
$\widetilde{\Delta R}=N_*^{-1}\sum_{c:n_c\ge2}(n_c-1)\widehat{\Delta R}_c$, a
multiplicative shrinkage of every cell's leave-one-out covariance gain by the factor
$(n_c-1)/n_c$.
\end{proposition}

\subsection{Inference}
\begin{theorem}[Asymptotic normality of the covariance-gain estimator]\label{thm:clt}
Suppose (i) $H_x(y)\in[-M,M]$ for all $x,y$; (ii) examples are partitioned into clusters
$g=1,\dots,G$ (e.g.\ images) that are mutually independent, with cluster sizes having a
finite third moment and $\max_g|g|=o(\sqrt{G})$; (iii) the limiting variance
$\sigma^2=\lim_{G\to\infty}\operatorname{Var}\!\bigl(\sqrt{N_*}(\widehat{\Delta R}-\theta)\bigr)$
exists and is strictly positive, where $\theta$ is the identified-subpopulation estimand
of Section~\ref{app:influence}; and (iv) $\phi$ takes values in a fixed, finite set,
each with strictly positive limiting population mass, so that every eligible cell's size
$n_c$ diverges almost surely as $N\to\infty$. Then
\begin{equation}
\sqrt{N_*}\bigl(\widehat{\Delta R}-\theta\bigr)\ \xrightarrow{d}\ \mathcal N(0,\sigma^2),
\end{equation}
and the plug-in sandwich variance $\widehat\sigma^2=N_*\dfrac{G}{G-1}\sum_{g}\bigl(\sum_{i\in g}\widehat\psi_i/N_*\bigr)^2$
of Section~\ref{app:influence} is consistent for $\sigma^2$.
\end{theorem}

\begin{corollary}[Relation to Diebold--Mariano/Giacomini--White]\label{cor:dm}
Under the trivial partition $\phi\equiv\text{const}$ (one cell containing every
example), $\widehat{\Delta T}=N^{-1}\sum_iH_i(y_i)$ is the sample mean of the paired loss
differential between the two models, and
Theorem~\ref{thm:clt} applied to $\widehat{\Delta T}$
with its image-clustered interval is the grouped, cluster-robust form of the
Diebold--Mariano statistic for that differential \citep{diebold1995comparing}; because
the conditioning information set is empty, it is equivalently a degenerate instance of
the Giacomini--White conditional predictive-ability test \citep{giacomini2006tests}.
\end{corollary}
